\subsection{Zweck der Identify-Phase}

Gemäss Handout ist die Identify-Phase die Orientierungsphase vor Discover.
Ziel ist es, die Ausgangslage mit dem \gls{zvv} gemeinsam zu klären, das Problemfeld zu strukturieren und ein gemeinsames Verständnis über Auftrag, Grenzen und Erkenntnisbedarf aufzubauen.

Die Phase dient damit als methodische Grundlage für den Einstieg in den ersten Diamanten (Discover/Define).
Nicht Lösungen entwickeln, sondern die Challenge präzisieren und die relevanten offenen Fragen für die Feldforschung ableiten.

\subsection{Eingesetzte Methoden und Vorgehen}

In der Identify-Phase werden vier aufeinander aufbauende Methoden eingesetzt.

\subsubsection{Challenge Deconstruction}
Die Design Challenge wird in ihre Kernbegriffe zerlegt (``Massnahmen'', ``vital'', ``autoabhängig'', ``Senior:innen'', ``Freizeitmobilität'').
Für jeden Begriff werden Assoziationen, implizite Annahmen und Spannungsfelder gesammelt und im Team verdichtet.

\subsubsection{Stakeholder Map}
Alle relevanten Akteur:innen werden in konzentrischen Kreisen erfasst: primäre Nutzende im Zentrum, danach Schlüssel-Partner und erweitertes Umfeld.
Die Methode wird genutzt, um Einflüsse auf Mobilitätsentscheidungen sichtbar zu machen und potenzielle Partner für spätere Interventionen zu identifizieren.

\subsubsection{Desk Research}
Als inhaltliche Basis werden bestehende Quellen zu Demografie, Mobilitätsverhalten, digitaler Nutzung und \gls{zvv}-Ausgangslage ausgewertet.
Verwendet werden insbesondere Bevölkerungsstatistik Kanton Zürich, Mikrozensus Mobilität und Verkehr, Mobilitätsdaten des Kantons Zürich, die Studie ``Digital Seniors'' sowie Unterlagen des \gls{zvv}.

\subsubsection{Design Space Map}
Die Ergebnisse werden abschliessend in einer Design Space Map synthetisiert.
Alle Erkenntnisse werden in drei Kategorien geordnet: ``Wir wissen'', ``Wir nehmen an'' und ``Wir wissen nicht''.
Diese Struktur dient als Brücke zur Discover-Phase und zur Schärfung des Interviewleitfadens.

\newpage
\subsection{Zentrale inhaltliche Ergebnisse}

\subsubsection{Auftrag, Zielgruppe und Rahmenbedingungen}
Der \gls{zvv} ist der grösste Verkehrsverbund der Schweiz und verfolgt drei strategische Ziele: höhere Auslastung in der \gls{nvz}, Verschiebung des Modalsplits zugunsten des \gls{oev} und nachhaltige Einstellungs- bzw. Verhaltensänderung.

Die bearbeitete Design Challenge lautet:

\begin{quote}
\textit{``Mit welchen Massnahmen können wir vitale und bislang überwiegend autoabhängige Senior:innen im Kanton Zürich zur verstärkten Nutzung des öffentlichen Verkehrs für Freizeitmobilität motivieren?''}
\end{quote}

Der Lösungsraum ist klar begrenzt.
Keine Eingriffe in Pricing, neue Tickets/Abos, Fahrplan oder Verkaufsstellen.
Im Fokus stehen Kommunikation, Services, Erlebnisse und persönliche Begleitung.

\subsubsection{Ergebnisse der Challenge Deconstruction}
Die Begriffsanalyse zeigt fünf zentrale Punkte:

\textbf{Massnahmen:} Erwartet werden keine tariflichen oder infrastrukturellen Eingriffe, sondern gestaltungsnahe Hebel im Service- und Kommunikationsraum.

\textbf{Vital:} Die Zielgruppe ist mehrheitlich selbständig, aktiv und konsumfreudig.
Gleichzeitig sehen sich viele nicht als ``Senior:innen'', was hohe Anforderungen an Sprache und Tonalität stellt.

\textbf{Autoabhängig:} Das Auto wird nicht nur als Transportmittel, sondern als Symbol für Kontrolle, Freiheit und Spontaneität verstanden.

\textbf{Senior:innen:} Keine homogene Gruppe.
Unterschiede zwischen aktiv/betagt sowie urban/ländlich und digitalaffin/offline sind für die Gestaltung zentral.

\textbf{Freizeitmobilität:} Der Nutzungskontext ist emotional und selbstgewählt.
Relevante Anlässe sind Ausflüge, Kultur, Sport, Kulinarik und Besuche.

\subsubsection{Zentrale Annahmen über die Zielgruppe}

Aus der Deconstruction verdichten sich vier Kernannahmen, die für die Gestaltung aller späteren Massnahmen zentral sind:

\textbf{Identität statt Alter:} Vitale Senior:innen sehen sich selbst oft nicht als Senioren.
Sie fühlen sich fit, aktiv und selbständig.
Diese psychologische Selbstwahrnehmung steht in Spannung zu ihrer chronologischen Zuordnung.
Folge: Sprache und Tonalität aller Kommunikation und Services muss diese Identität respektieren.
``Senioren-Angebote'' mit paternalistischem Unterton funktionieren nicht; es braucht statt dessen eine Ansprache als selbstbestimmte, erlebnisorientierte Personen.

\textbf{Autoabhängigkeit als psychologisches Phänomen:} Das Auto ist nicht bloss ein Transportmittel, sondern ein psychologisches Symbol.
Es steht für Kontrolle über Zeit, Tempo und Ziel, für Freiheit und Autonomie, für das Festhalten an Unabhängigkeit.
Diese emotionale Bindung erklärt, warum Preis- oder Angebotargumente oft nicht greifen.
Es braucht statt dessen Angebote, die das gleiche psychologische Bedürfnis (Kontrolle, Freiheit, Selbständigkeit) mit anderen Mitteln erfüllen.

\textbf{Drei unterschiedliche Segmente:} Die Zielgruppe ist heterogen nach drei Dimensionen: Lebensstil (aktiv vs. betagte), Geografie (urban vs. ländlich) und Digitalität (online vs. offline).
Diese drei Segmente haben unterschiedliche Lebenswelten, unterschiedliche Vertrauenspersonen, unterschiedliche Informationskanäle und unterschiedliche Barrieren.
Damit ist klar, dass ein ``One-Size-Fits-All''-Ansatz nicht funktioniert.

\textbf{Freizeitfahrten als Genuss und Erlebnis:} Anders als Alltagsfahrten (zur Arbeit, zum Einkaufen) sind Freizeitfahrten selbstgewählt und emotional besetzt.
Im Vordergrund stehen Genuss, Erlebnis und soziale Begegnung, nicht Effizienz oder Preis.
Das bedeutet: Ein gutes Freizeitangebot muss nicht bloss mobil versorgen, sondern eine sinnliche und soziale Erfahrung bieten.
Der Weg ist Teil des Erlebnisses, nicht bloss Mittel zum Zweck.

\subsubsection{Visualisierung der Challenge Deconstruction}
Die Mindmap basiert auf der von der Arbeitsgruppe erstellten Deconstruction-Mindmap und wird als skalierbare Vektorgrafik in LaTeX nachgebaut.
Sie verdichtet die zentralen Begriffe der Design Challenge und macht Zusammenhänge, Annahmen und Spannungsfelder auf einen Blick sichtbar.
\newpage

\FullWidthFigure[\textwidth]{sections/figures/identify_deconstruction_mindmap}{Challenge Deconstruction im \gls{zvv}-Praxis-Case (eigene Darstellung auf Basis der Gruppen-Mindmap)}

\subsubsection{Ergebnisse der Stakeholder Map}

Die Stakeholder Map erfasst alle relevanten Akteur:innen in konzentrischen Kreisen - vom Zentrum mit der Zielgruppe bis zum gesellschaftlichen Umfeld.

\textbf{Primäre Nutzende (Zentrum):}
Im innersten Kreis stehen zwei zentrale Gruppen:
Vitale, autoabhängige Senior:innen (65--79 Jahre) als Hauptzielgruppe: mobil, kaufkräftig, gewohnheitsgeprägt.
Bereits öV-fahrende Senior:innen als Referenzgruppe: wertvoll für Erfahrungsaustausch und Peer-Motivation.
Zusätzlich erfasst die Map das soziale Umfeld (Partner:in, Familie, Freunde, Kinder) als entscheidende Einflussgrösse.
Ein erstes Kernergebnis deutet sich ab: Der Mobilitätsentscheid ist oft nicht individuell - das Umfeld beeinflusst die Wahl stark.

\textbf{Schlüssel-Partner (mittlerer Kreis):}
Die Stakeholder Map identifiziert diverse Partnerorganisationen:
\gls{zvv} (Auftraggeber): Koordination, Kommunikation, Marke - aber kein direktes Fahrerlebnis.
Verkehrsunternehmen (\gls{sbb}, \gls{zsg}, weitere Verbunde): Das Fahrerlebnis entsteht hier, nicht beim \gls{zvv}.
Pro Senectute und Verein Zürcher Seniorinnen und Senioren: grösste Seniorenorganisationen mit existierendem Vertrauen und direktem Zugang.
Gemeinden und Quartierzentren: niederschwelliger physischer Zugang, besonders für Offline-Senior:innen.
Alltagsnahe Kontaktpunkte (Apotheken, Bibliotheken, Hausärzt:innen): Vertrauenspersonen im Alltag, geeignet für analoge Infovermittlung.
Freizeitanbieter: Kultur, Sport, Kulinarik, Ausflugsziele als potenzielle Kooperationspartner.
``Invisible Influencers'' (Ärzt:innen, Gemeinden, Versicherungen, Medien): prägen das gesellschaftliche Bild von Senior:innen und Mobilität, ohne selbst Angebote zu schaffen.

\textbf{Weiteres Umfeld (äusserer Kreis):}
Bund und Kantone (Regulierung, Förderung, Mobilitätspolitik).
Tourismusorganisationen (Freizeitziele und Ausflugspakete als Motivationsanker).
Autolobby (Gegenspieler im Narrativ).
Gesellschaft Schweiz (Normen rund um Mobilität, Selbständigkeit und Alter).

\textbf{Zentrale Erkenntnisse:}
Kernergebnis: Der \gls{zvv} erreicht diese Zielgruppe vermutlich nicht allein.
Partnerorganisationen mit bestehendem Vertrauen spielen eine entscheidende Rolle.
Das soziale Umfeld - Partner:in, Familie, Freunde - hat einen grösseren Einfluss auf den Mobilitätsentscheid, als die Fragestellung zunächst vermuten liess.

\FullWidthFigure[\textwidth]{sections/figures/stakeholder_map}{Stakeholder Map des \gls{zvv}-Praxis-Case (eigene Darstellung)}

\subsubsection{Ergebnisse der Desk Research}

\textbf{Methodisches Vorgehen:}
Wir sichteten die wichtigsten Schweizer Quellen zur Zielgruppe mit dem Ziel eines datengestützten Bildes: Mobilitätsverhalten, demografische Ausgangslage und digitale Nutzung - differenziert nach Segmenten, wo Daten verfügbar waren.

\textbf{Demografische Ausgangslage:}
Im Kanton Zürich leben rund 327'000 Personen über 65 Jahre (21\% der Gesamtbevölkerung), mit erwarteter Zunahme auf ca. 420'000 bis 2040 (+28\%).
Das grösste Hebelsegment sind aktive Senior:innen (65--79 Jahre): ca. 190'000 Personen, kaufkräftig mit mittlerem Nettovermögen von CHF 420'000.
Betagte Senior:innen (75--85 Jahre): ca. 107'000 Personen in Mobilitätsübergangsphase.
Offline-Senior:innen (kein Internet/Smartphone): ca. 30'000 Personen, ausschliesslich auf analoge Kanäle angewiesen.

\textbf{Mobilitätsverhalten nach Segment:}
Aktive Senior:innen (65--75 Jahre): durchschnittlich 27,4 km pro Tag - 57\% mit Auto, nur 18\% öV (CH-Ø: 23\%, Kanton Zürich ca. 22\%).
Über 60\% der Freizeitfahrten mit Auto, auch im städtischen Raum.
78\% aller Ü65 sind in der Nebenverkehrszeit unterwegs (ideal für \gls{zvv}).
Hauptmotive für Autonutzung: Flexibilität (38\%), Bequemlichkeit (27\%), kein Gepäckproblem (21\%), Sicherheitsgefühl (14\%).

Betagte Senior:innen (75--85 Jahre): durchschnittlich nur noch 11 km pro Tag - Mobilitätsradius schrumpft.
MIV-Anteil ca. 55\%, öV-Anteil steigt auf ca. 22\%.
Sicherheitsbedenken im öV: Orientierungslosigkeit, Umstiegsunsicherheit, fehlende Assistenz.

Stadt-Land-Gefälle: Stadt Zürich öV-Anteil Ü65 ca. 28\%, ländliche Gemeinden nur ca. 11\%.
Im ländlichen Raum oft stündlicher Takt und fehlende Direktverbindungen zu Freizielzielen.

\textbf{Digitale Nutzung nach Segment:}
Aktive 65--74 Jahre: 84\% Internetnutzung, 74\% Smartphone - digitale Kanäle (App, Online-Ticket) gut erreichbar.
Betagte 75--84 Jahre: 62\% Internet, 49\% Smartphone - Gap wächst, analoge Ergänzung nötig.
Offline-Senior:innen (Ø 78 J.): unter 25\% - ausschliesslich analoge Zugänge.
Hauptgründe für Nicht-Nutzung: ``zu kompliziert'' (41\%), Sicherheitsbedenken (29\%), hoher Lernaufwand (22\%).

\textbf{Der kritische Moment: Führerscheinabgabe:}
Ca. 190'000 Personen Ü65 im Kanton Zürich besitzen einen Führerausweis (ca. 80\% der Zielgruppe).
Ca. 8'000--10'000 Personen geben den Führerausweis pro Jahr ab - ein starkes Aktivierungsfenster für den öV.
Nur 34\% bereiten sich im Vorfeld mit öV-Alternativen vor.
Hauptbarrieren nach Abgabe: Orientierungslosigkeit (43\%), Sicherheitsbedenken (28\%), fehlende Information (19\%).
Bei früher Vorbereitung (mindestens 1 Jahr vorher) ist die spätere öV-Nutzungsrate signifikant höher.

\textbf{Nutzungsbarrieren für den öV:}
Psychologisch: Auto = Freiheit, Kontrolle, Spontanität; öV = Abhängigkeit. Gewohnheit über Jahrzehnte. Erste negative Erfahrung wirkt dauerhaft abschreckend. Soziale Norm: Auto dominiert im Umfeld.

Praktisch: Gepäck und Einkäufe im Auto bequemer. Wege zur Haltestelle bei Regen/Kälte als Hürde. ``Letzter Kilometer'': Freizeitziele oft nicht direkt erreichbar. öV erfordert Planung, Auto ist spontan verfügbar.

Systemisch: Tarifkomplexität, digitaler Ausschluss (Apps/Online-Tickets), abbauende physische Schalter, Stadt-Land-Gefälle bei Direktverbindungen, wenig sichtbare \gls{zvv}-Kommunikation für diese Zielgruppe.

\subsubsection{Ergebnisse der Design Space Map}

Abschliessend verdichteten wir alle Erkenntnisse in einer Design Space Map - ein dreigeteiltes Raster, das den Wissensstand systematisch offenlegt.

\textbf{Was wir wissen (Rechte Spalte):}
Züge sind nicht ausgelastet zwischen Pendlerverkehr und Nebenverkehrszeit - \gls{nvz}-Potenzial vorhanden.
Das Auto ist das Hauptfortbewegungsmittel der Zielgruppe.
\gls{zvv} hat bereits Angebote, erreicht die Zielgruppe damit aber nicht.
Die Zielgruppe ist kaufkräftig mit hohem Nettovermögen.
Die Kosten des Autos werden von der Zielgruppe oft nicht vollständig wahrgenommen.
\gls{zvv} kennt seine Kund:innen in dieser Altersgruppe kaum im Detail.
Urban gibt es guten öV-Zugang, ländlich ist der öV-Zugang deutlich schlechter.
Die Zielgruppe ist digitalaffin bei 65--74 Jahren, die Affinität nimmt ab 75 ab.

\textbf{Was wir annehmen (Mittlere Spalte):}
Das Auto steht für Freiheit, Spontanität und Selbständigkeit - nicht bloss als Transportmittel.
Physische Barrieren hemmen die öV-Nutzung: Gepäck, Wetter, schwer erreichbare Ziele (ländlich vs. Stadt).
Freizeitfahrten sind emotional besetzt und selbstgewählt, nicht zweckorientiert.
Gewohnheit schlägt rationale Entscheidung - das Auto ist die Standardwahl.
Die erste negative öV-Erfahrung (Ticket, Umstieg, Verspätung) wirkt dauerhaft abschreckend.
Social Influence ist stark: Partner:in, Freunde und Familie beeinflussen den Mobilitätsentscheid erheblich.
Die Zielgruppe sieht sich nicht als ``Senioren'' - sie fühlt sich fit und selbständig.
Beim öV entsteht Unsicherheit bei Ticket, Umstieg und Abhängigkeit vom System.

\textbf{Was wir nicht wissen (Linke Spalte):}
Gute und schlechte öV-Erfahrungen der Zielgruppe im Detail.
Konkrete Verhaltensgründe, weshalb Auto dem öV vorgezogen wird.
Typischer Tagesablauf und Freizeitprioritäten.
Informationsquellen der Zielgruppe.
Wo sie ihre Schwerpunkte setzen - die Highlight-Momente des Tages.
Wann wird überhaupt über Verkehrsmittelwahl entschieden?
Was hält sie konkret vom öV ab - was könnte sie motivieren?
Freizeitverhalten und digitale Affinität im Detail.

Mit diesem Wissensstand starteten wir die Discover-Phase.
Die offenen Fragen und Annahmen bildeten die Grundlage für den Interviewleitfaden.

\subsubsection{Visualisierung der Design Space Map}
Die Design Space Map ordnet den gesamten Wissensstand in einem dreigeteilten Raster:

\textbf{Linke Spalte -- Was wir nicht wissen:} Hier sind alle Fragen gebündelt, die für eine Aktivierungsstrategie zentral sind, aber noch offene Forschungsfragen darstellen.
Dazu gehören konkrete Gründe für die Autoabhängigkeit, Informationsquellen der Zielgruppe, situative Entscheidungslogiken, digitale Affinität und die alltägliche Routinegestaltung.

\textbf{Mittlere Spalte -- Was wir annehmen:} Basierend auf Desk Research und Teamdiskurs werden hier hypothetische Erklärungen festgehalten.
Sie beziehen sich auf psychologische Faktoren (Kontrollgefühl, Sicherheit), systemische Barriereeffekte (Erreichbarkeit, Infrastruktur) sowie soziale Einflussfaktoren.
Diese Annahmen sind bewusst gekennzeichnet und werden in der Discover-Phase empirisch überprüft.

\textbf{Rechte Spalte -- Was wir wissen:} Hier sind die gesicherten Erkenntnisse aus Sekundärdaten und Behördenkommunikation aufgeführt: die Zielgruppendemografia, die Verkehrsmittelpräferenzen, das Potenzial in der \gls{nvz} sowie die bestehende Angebotspalette des \gls{zvv}.

Die Map dient als Kompass für die kommende Feldforschung.
Sie zeigt, welche Fragen mit Interview und Beobachtung geklärt werden müssen, um von Annahmen zu verifizierten Insights zu gelangen.

\FullWidthFigure[\textwidth]{sections/figures/design_space_map}{Design Space Map des \gls{zvv}-Praxis-Case (eigene Darstellung)}

\subsection{Reflexion der Methode}

Die Kombination aus Challenge Deconstruction, Stakeholder Map, Desk Research und Design Space Map erweist sich als sinnvoll, um vom breiten Problemraum zu einem klar strukturierten Forschungsauftrag zu gelangen.
Insbesondere die Trennung in ``wissen/nehmen an/wissen nicht'' hilft, Annahmen sichtbar zu machen und nicht vorschnell als Fakten zu behandeln.

Theoretisch lässt sich dieses Vorgehen im Kontext des ``Fuzzy Front End'' verorten, das \textcite{lewrick2018} als entscheidende Weichenstellung im Design-Thinking-Prozess beschreiben.
Bevor divergent geforscht wird, muss der Problemraum eingegrenzt und das Team auf ein gemeinsames Verständnis ausgerichtet werden.
Die Design Space Map erfüllt genau diese Funktion, indem sie den Übergang vom impliziten Teamwissen zu einer expliziten, teilbaren Wissensstruktur formalisiert.

Gleichzeitig zeigt sich eine methodische Grenze der Identify-Phase.
Sie schafft Orientierung, liefert aber keine belastbare Evidenz zu Motiven und situativen Entscheidungen aus erster Hand.
Genau daraus ergibt sich der nächste Schritt im Prozess: qualitative Interviews in der Discover-Phase, um die offenen Fragen empirisch zu klären.

Rückblickend passen wir zwei Aspekte an.
Erstens versehen wir die Stakeholder Map stärker mit Hypothesen über Einflussrichtungen (z.\,B. ``Partner:in beeinflusst Verkehrsmittelwahl positiv/negativ''), um in den Interviews gezielter nachfragen zu können.
Zweitens ist der Desk Research zeitlich knapp bemessen.
Eine systematischere Literaturrecherche zu Verhaltensänderungsmodellen (etwa Transtheoretisches Modell oder Theory of Planned Behaviour) schärft die Annahmen in der Design Space Map bereits früher theoretisch.
